Powerful Number（以下简称 PN）筛类似于杜教筛，或者说是杜教筛的一个扩展，可以拿来求一些积性函数的前缀和。

\textbf{要求}：
\begin{itemize}
    \item 存在一个函数 $g$ 满足：
    \begin{itemize}
        \item $g$ 是积性函数；
        \item $g$ 易求前缀和；
        \item 对于质数 $p$，$g(p) = f(p)$。
    \end{itemize}
\end{itemize}

假设现在要求积性函数 $f$ 的前缀和 $F(n) = \sum_{i=1}^{n} f(i)$。

\subsubsection{Powerful Number}

\textbf{定义}：对于正整数 $n$，记 $n$ 的质因数分解为 $n = \prod_{i=1}^{m} p_i^{e_i}$。$n$ 是 PN 当且仅当 $\forall 1 \le i \le m,\ e_i > 1$。

\textbf{性质 1}：所有 PN 都可以表示成 $a^{2}b^{3}$ 的形式。

\textbf{证明}：若 $e_i$ 是偶数，则将 $p_i^{e_i}$ 合并进 $a^{2}$ 里；若 $e_i$ 为奇数，则先将 $p_i^{3}$ 合并进 $b^{3}$ 里，再将 $p_i^{e_i-3}$ 合并进 $a^{2}$ 里。

\textbf{性质 2}：$n$ 以内的 PN 至多有 $O(\sqrt{n})$ 个。

\textbf{证明}：考虑枚举 $a$，再考虑满足条件的 $b$ 的个数，有 PN 的个数约等于
\[
\int_{1}^{\sqrt{n}} \sqrt[3]{\frac{n}{x^2}} \,\mathrm{d}x = O(\sqrt{n})。
\]

那么如何求出 $n$ 以内所有的 PN 呢？线性筛找出 $\sqrt{n}$ 内的所有素数，再 DFS 搜索各素数的指数即可。由于 $n$ 以内的 PN 至多有 $O(\sqrt{n})$ 个，所以至多搜索 $O(\sqrt{n})$ 次。

\subsubsection{PN 筛}

首先，构造出一个易求前缀和的积性函数 $g$，且满足对于素数 $p$，$g(p) = f(p)$。记 $G(n) = \sum_{i=1}^{n} g(i)$。

然后，构造函数 $h = f / g$，这里的 $/$ 表示狄利克雷卷积除法。根据狄利克雷卷积的性质可以得知 $h$ 也为积性函数，因此 $h(1) = 1$。$f = g * h$，这里 $*$ 表示狄利克雷卷积。

对于素数 $p$，
\[
f(p) = g(1)h(p) + g(p)h(1) = h(p) + g(p) \implies h(p) = 0。
\]
根据 $h(p)=0$ 和 $h$ 是积性函数可以推出对于非 PN 的数 $n$ 有 $h(n) = 0$，即 $h$ 仅在 PN 处取有效值。

现在，根据 $f = g * h$ 有
\[
\begin{aligned}
F(n) &= \sum_{i = 1}^{n} f(i) \\
     &= \sum_{i = 1}^{n} \sum_{d\mid i} h(d) g\!\left(\frac{i}{d}\right) \\
     &= \sum_{d=1}^{n} \sum_{i=1}^{\lfloor n/d \rfloor} h(d) g(i) \\
     &= \sum_{d=1}^{n} h(d) \sum_{i=1}^{\lfloor n/d \rfloor} g(i) \\
     &= \sum_{d=1}^{n} h(d) G\!\left(\left\lfloor \frac{n}{d} \right\rfloor\right) \\
     &= \sum_{\substack{d=1 \\ d \text{ 是 PN}}}^{n} h(d) G\!\left(\left\lfloor \frac{n}{d} \right\rfloor\right)。
\end{aligned}
\]

$O(\sqrt{n})$ 找出所有 PN，计算出所有 $h$ 的有效值。对于 $h$ 有效值的计算，只需要计算出所有 $h(p^c)$ 处的值，就可以根据 $h$ 为积性函数推出 $h$ 的所有有效值。现在对于每一个有效值 $d$，计算 $h(d) G\bigl(\lfloor n/d \rfloor\bigr)$ 并累加即可得到 $F(n)$。

下面考虑计算 $h(p^c)$，一共有两种方法：一种是直接推出 $h(p^c)$ 仅与 $p, c$ 有关的计算公式，再根据公式计算 $h(p^c)$；另一种是根据 $f = g * h$ 有 $f(p^c) = \sum_{i=0}^{c} g(p^i) h(p^{c-i})$，移项可得
\[
h(p^c) = f(p^c) - \sum_{i=1}^{c} g(p^i) h(p^{c-i})，
\]
现在就可以枚举素数 $p$ 再枚举指数 $c$ 求解出所有 $h(p^c)$。

\subsubsubsection{过程}
\begin{enumerate}
    \item 构造 $g$；
    \item 构造快速计算 $G$ 的方法；
    \item 计算 $h(p^c)$；
    \item 搜索 PN，过程中累加答案；
    \item 得到结果。
\end{enumerate}
对于第 3 步，可以直接根据公式计算，可以使用枚举法预处理打表，也可以搜索到了再临时推。

\subsubsubsection{性质}
以使用第二种方法计算 $h(p^c)$ 为例进行分析。可以分为计算 $h(p^c)$ 和搜索两部分进行分析。

对于第一部分，根据 $\sqrt{n}$ 内的素数个数为 $O\!\left(\frac{\sqrt{n}}{\log n}\right)$，每个素数 $p$ 的指数 $c$ 至多为 $\log n$，计算 $h(p^c)$ 需要循环 $(c-1)$ 次，由此有第一部分的时间复杂度为 $O\!\left(\frac{\sqrt{n}}{\log n} \cdot \log n \cdot \log n\right) = O(\sqrt{n}\log n)$，且这是一个宽松的上界。根据题目的不同还可以添加不同的优化，从而降低第一部分的时间复杂度。

对于搜索部分，由于 $n$ 以内的 PN 至多有 $O(\sqrt{n})$ 个，所以至多搜索 $O(\sqrt{n})$ 次。对于每一个 PN，根据计算 $G$ 的方法不同，时间复杂度也不同。例如，假设计算 $G\!\left(\lfloor n/d \rfloor\right)$ 的时间复杂度为 $O(1)$，则第二部分的复杂度为 $O(\sqrt{n})$。

特别地，若借助杜教筛计算 $G\!\left(\lfloor n/d \rfloor\right)$，则第二部分的时间复杂度为杜教筛的时间复杂度，即 $O(n^{2/3})$。因为若事先计算一次 $G(n)$，并且预先使用线性筛优化和用支持快速随机访问的数据结构（如 C++ 中的 `std::map` 和 `std::unordered\_map`）记录较大的值，则杜教筛过程中用到的 $G\!\left(\lfloor n/d \rfloor\right)$ 都是线性筛中记录的或者 `std::map`' 中记录的，这一点可以直接用程序验证。

对于空间复杂度，其瓶颈在于存储 $h(p^c)$。若使用二维数组 $a$ 记录，$a_{i,j}$ 表示 $h(p_i^j)$ 的值，则空间复杂度为 $O\!\left(\frac{\sqrt{n}}{\log n} \cdot \log n\right) = O(\sqrt{n})$。